\section*{Bab \ref{ch:b1:periodicity} --- Keperiodikan}

\begin{solution}{exo:b1:periodicity:1}
\ce{Ca} $[\ce{Ar}]\,4\text{s}^2$: \omterm{def:b1:periodicity:table}{periode} 4, \omterm{def:b1:periodicity:table}{golongan} 2, \omterm{def:b1:periodicity:table}{blok} s. \ce{Fe}
$[\ce{Ar}]\,3\text{d}^6 4\text{s}^2$: \omterm{def:b1:periodicity:table}{periode} 4, \omterm{def:b1:periodicity:table}{golongan} 8, \omterm{def:b1:periodicity:table}{blok} d.
\ce{Se} $[\ce{Ar}]\,3\text{d}^{10} 4\text{s}^2 4\text{p}^4$: \omterm{def:b1:periodicity:table}{periode} 4,
\omterm{def:b1:periodicity:table}{golongan} 16, \omterm{def:b1:periodicity:table}{blok} p. \ce{I} $[\ce{Kr}]\,4\text{d}^{10} 5\text{s}^2 5\text{p}^5$: \omterm{def:b1:periodicity:table}{periode} 5, \omterm{def:b1:periodicity:table}{golongan} 17, \omterm{def:b1:periodicity:table}{blok} p.
\end{solution}

\begin{solution}{exo:b1:periodicity:2}
\ce{Na} lebih besar daripada \ce{Cl}: seperiode, $Z^{*}$ lebih kecil.
\ce{Na} lebih besar daripada \ce{Na+}, yang telah kehilangan kulit
3s-nya. \ce{Cl-} (\qty{181}{pm}) lebih besar daripada \ce{F-}
(\qty{133}{pm}): satu kulit lebih banyak. \ce{Na+} (\qty{102}{pm}) lebih
besar daripada \ce{Mg^{2+}} (\qty{72}{pm}): sepuluh elektron yang sama,
muatan inti 11 lawan 12.
% ledger: rion:Cl-VI, rion:F-VI, rion:Na+VI, rion:Mg2+VI
\end{solution}

\begin{solution}{exo:b1:periodicity:3}
$\ce{Na} (5.14) < \ce{Al} (5.99) < \ce{Mg} (7.65) < \ce{S} (10.36) <
\ce{P} (10.49) < \ce{Ar} (15.76)$. Kenaikan umum sepanjang \omterm{def:b1:periodicity:table}{periode}
terputus dua kali: \ce{Al} di bawah \ce{Mg} (elektron yang dilepaskan
adalah elektron 3p pertama, yang energinya lebih tinggi daripada 3s) dan
\ce{S} di bawah \ce{P} (elektron yang dilepaskan adalah elektron 3p
pertama yang berpasangan).
\end{solution}

\begin{solution}{exo:b1:periodicity:4}
\omterm{def:b1:periodicity:polarisability}{Polarisabilitas} mengukur seberapa mudah awan elektron diubah bentuknya
oleh medan listrik. \ce{I-} lebih mudah terpolarisasi daripada \ce{F-}:
ion itu lebih besar dan elektron-elektron terluarnya lebih jauh dari inti.
\ce{Cl-} (18 elektron yang diikat oleh $Z = 17$) jauh lebih mudah
terpolarisasi daripada \ce{Na+} (10 elektron yang diikat oleh $Z = 11$),
seperti halnya anion dibandingkan dengan kation.
\end{solution}

\begin{solution}{exo:b1:periodicity:5}
Rasio-rasionya: $18.83/5.99 = 3.1$, $28.45/18.83 = 1.5$,
$119.99/28.45 = 4.2$. Lonjakan terjadi sesudah elektron ketiga: tiga
\omterm{def:b1:quantum-numbers:configuration}{elektron valensi}, \omterm{def:b1:periodicity:table}{golongan} 13, \omterm{def:b1:periodicity:table}{periode} 3: aluminium. Unsur ini membentuk
\ce{Al^{3+}}, $1\text{s}^2 2\text{s}^2
2\text{p}^6$.
% ledger: ie:Al, ie2:Al, ie3:Al, ie4:Al
\end{solution}

\begin{solution}{exo:b1:periodicity:6}
Elektron tambahan \ce{F-} memasuki kulit 2p yang kecil, tempat elektron
itu ditolak kuat oleh tujuh \omterm{def:b1:quantum-numbers:configuration}{elektron valensi} yang sudah ada; di dalam
kulit 3p klor yang lebih besar tolakannya lebih kecil, sehingga energi
yang dilepaskan lebih banyak. Pada nitrogen ($2\text{p}^3$,
\omorbs{u,u,u}) elektron tambahan harus berpasangan di dalam \omterm{def:b1:quantum-numbers:orbital}{subkulit}
setengah penuh; pada magnesium ($3\text{s}^2$) elektron itu harus
memulai \omterm{def:b1:quantum-numbers:orbital}{subkulit} 3p yang lebih tinggi. Dalam kedua kasus anionnya tidak
lebih stabil daripada atom beserta elektron bebas.
\end{solution}

\begin{solution}{exo:b1:periodicity:7}
$\chi_M(\ce{Na}) = (5.14 + 0.55)/2 = \qty{2.85}{eV}$; $\chi_M(\ce{Cl}) =
(12.97 + 3.61)/2 = \qty{8.29}{eV}$. Selisihnya besar: pasangan elektron
ikatan \ce{Na-Cl} hampir sepenuhnya milik klor, dan natrium klorida
bersifat ionik, \ce{Na+ Cl-}.
\end{solution}

\begin{solution}{exo:b1:periodicity:8}
\ce{H-F}: $\Delta\chi = 1.78$, di perbatasan daerah ionik (dalam praktik
ikatan kovalen yang sangat polar), \ce{F^{\delta-}}. \ce{C-H}: 0.35,
hampir nonpolar. \ce{Na-Cl}: 2.23, ionik, \ce{Cl^{-}}. \ce{C-Cl}: 0.61,
kovalen polar, \ce{C^{\delta+}-Cl^{\delta-}}. \ce{O-H}: 1.24, kovalen
polar, \ce{O^{\delta-}-H^{\delta+}}.
\end{solution}

\begin{solution}{exo:b1:periodicity:9}
Menuruni \omterm{def:b1:periodicity:table}{golongan}, \omterm{def:b1:quantum-numbers:configuration}{elektron valensi} berada pada kulit dengan $n$ yang
lebih besar, lebih jauh dari inti dan ditapis oleh lebih banyak \omterm{def:b1:quantum-numbers:configuration}{elektron
teras}; elektron itu lebih mudah dilepaskan. Rubidium, satu kulit lebih
jauh, mempunyai \omterm{def:b1:periodicity:ionisation-energy}{energi ionisasi pertama} di bawah \qty{4.34}{eV}.
\end{solution}

\begin{solution}{exo:b1:periodicity:10}
Seng adalah $[\ce{Ar}]\,3\text{d}^{10} 4\text{s}^2$: elektronnya berasal
dari 4s yang penuh; galium, $\dots 4\text{s}^2 4\text{p}^1$, melepaskan
satu-satunya elektron 4p, yang energinya lebih tinggi dan ditapis oleh 4s
dan 3d yang penuh. Arsen, $4\text{p}^3$, melepaskan sebuah elektron dari
\omterm{def:b1:quantum-numbers:orbital}{subkulit} setengah penuh; selenium, $4\text{p}^4$, melepaskan elektron
yang berpasangan, yang terdorong naik oleh tolakan pasangannya: energi
ionisasinya sedikit lebih rendah.
\end{solution}

\begin{solution}{exo:b1:periodicity:11}
$\Delta E = (3.98 - 2.20)^2 = 1.78^2 = \qty{3.17}{eV} = 3.17 \times 96.5 =
\qty{306}{kJ/mol}$. Maka $D(\ce{H-F}) = \tfrac12(436 + 158) + 306 = 297 +
306 \approx \qty{603}{kJ/mol}$. Kelebihan yang besar itu adalah tarikan
tambahan antara muatan parsial $\ce{H^{\delta+}}$ dan $\ce{F^{\delta-}}$,
akibat selisih \omterm{def:b1:periodicity:electronegativity}{keelektronegatifan} yang besar.
\end{solution}

\begin{solution}{exo:b1:periodicity:12}
\omterm{def:b1:periodicity:electronegativity}{Skala Pauling} didefinisikan dari energi ikatan A--B: unsur yang tidak
membentuk ikatan tidak mempunyai nilai Pauling. Helium, neon, dan argon
tidak membentuk senyawa yang stabil. Kripton dan terutama xenon mempunyai
kulit valensi yang lebih besar dan lebih mudah terpolarisasi serta energi
ionisasi yang lebih rendah (\qty{14.0}{eV} untuk kripton dibandingkan
dengan \qty{15.76}{eV} untuk argon); keduanya dioksidasi oleh fluorin dan
oksigen dan membentuk senyawa seperti \ce{XeF2}, sehingga energi
ikatannya dapat diukur.
% ledger: ie:Kr, ie:Ar
\end{solution}

\begin{solution}{pb:b1:periodicity:1}
\textbf{1.} Setiap elektron dilepaskan dari ion yang lebih positif dan
lebih kecil daripada ion sebelumnya, yang mengikat elektronnya lebih kuat.
\textbf{2.} $15.04/7.65 = 1.97$; $80.14/15.04 = 5.33$; $109.27/80.14 = 1.36$.
\textbf{3.} Lonjakan terjadi sesudah dua elektron: X mempunyai dua
\omterm{def:b1:quantum-numbers:configuration}{elektron valensi}, \omterm{def:b1:periodicity:table}{golongan} 2; pada \omterm{def:b1:periodicity:table}{periode} 3 unsur itu adalah magnesium.
\textbf{4.} \ce{Mg^{2+}}, $1\text{s}^2 2\text{s}^2 2\text{p}^6$,
konfigurasi neon.
\textbf{5.} Logam: \omterm{def:b1:periodicity:table}{golongan} 2, di sebelah kiri tabel, dengan \omterm{def:b1:periodicity:ionisation-energy}{energi
ionisasi pertama} yang rendah.
\textbf{6.} Elektron ketiga berasal dari teras 2p, dengan $n$ yang lebih
kecil, jauh lebih dekat ke inti dan hampir tidak tertapis, dan dilepaskan
dari ion bermuatan dua.
% ledger: ie:Mg, ie2:Mg, ie3:Mg, ie4:Mg
\textbf{7.} Semuanya mempunyai sepuluh elektron: $1\text{s}^2 2\text{s}^2 2\text{p}^6$.
\textbf{8.} Sepuluh elektron yang sama diikat oleh muatan inti 8, 9, 11,
12, 13: awannya menyusut seiring bertambahnya $Z$.
\textbf{9.} \ce{O^{2-}}: yang terbesar, dengan muatan inti terendah untuk
kesepuluh elektronnya, dan berupa anion.
\textbf{10.} $140/54 = 2.6$.
\textbf{11.} $198.8/2 = \qty{99.4}{pm}$; anionnya $181/99.4 = 1.8$ kali
lebih besar daripada \omterm{def:b1:periodicity:radius}{jari-jari kovalennya}.
% ledger: re:Cl2, rion:Cl-VI
\textbf{12.} $\ce{Mg^{2+}} < \ce{Na+} < \ce{Cl-}$: kedua kation
isoelektronik (muatan 12 lawan 11), dan \ce{Cl-} mempunyai kulit ketiga.
\textbf{13.} $\chi_M$ = 10.41 (\ce{F}), 8.29 (\ce{Cl}), 7.59 (\ce{Br}),
7.53 (\ce{O}), 6.26 (\ce{C}), dalam eV.
\textbf{14.} Mulliken: $\ce{F} > \ce{Cl} > \ce{Br} > \ce{O} > \ce{C}$;
Pauling: $\ce{F} > \ce{O} > \ce{Cl} > \ce{Br} > \ce{C}$. Oksigen naik
dari urutan keempat ke urutan kedua.
\textbf{15.} $3.98/10.41 = 0.382$, $3.16/8.29 = 0.381$, $2.96/7.59 =
0.390$: hampir tetap, sekitar 0.38; untuk halogen kedua skala itu
sebanding.
\textbf{16.} \omterm{def:b1:periodicity:electron-affinity}{Afinitas elektron} oksigen kecil (\qty{1.44}{eV}) karena
elektron tambahan harus berpasangan di dalam kulit 2p yang kompak; nilai
Mulliken atom bebas meremehkan tarikan yang dikerahkan oksigen di dalam
ikatan-ikatannya, yang diukur langsung oleh \omterm{def:b1:periodicity:electronegativity}{skala Pauling}, yang disusun
dari ikatan.
\textbf{17.} Klor dalam \ce{HCl}; fluorin dalam \ce{ClF}.
\textbf{18.} $\Delta\chi = 0.96$: kovalen polar.
\textbf{19.} $D(\ce{H-H}) = 2 \times 218.0 = \qty{436.0}{kJ/mol}$.
\textbf{20.} $D(\ce{Cl-Cl}) = 2 \times 121.3 = \qty{242.6}{kJ/mol}$.
\textbf{21.} $D(\ce{H-Cl}) = 218.0 + 121.3 - (-92.3) = \qty{431.6}{kJ/mol}$.
\textbf{22.} $\Delta E = 431.6 - \tfrac12(436.0 + 242.6) = 431.6 - 339.3 =
\qty{92.3}{kJ/mol}$. Dalam lambang, $D(\ce{H-Cl}) = \Delta_fH(\ce{H}) +
\Delta_fH(\ce{Cl}) - \Delta_fH(\ce{HCl})$, sedangkan rata-rata ikatan
simetrisnya adalah $\Delta_fH(\ce{H}) + \Delta_fH(\ce{Cl})$, sehingga
$\Delta E =
-\Delta_fH(\ce{HCl})$.
\textbf{23.} $92.3/96.485 = \qty{0.957}{eV}$.
\textbf{24.} $|\chi(\ce{Cl}) - \chi(\ce{H})| = \sqrt{0.957} =
\textbf{0.98}$, dibandingkan dengan $3.16 - 2.20 = 0.96$ dalam tabel:
aritmetika Pauling, yang dikerjakan dengan data modern, menghasilkan
selisih dalam tabel dengan ketelitian 0.02.
\end{solution}
